% defined-in: spivak (page 236)
% entity-id: prop-rolles-zero-derivative
% entity-type: prop
% macro: \RollesZeroDerivative
% args: 1
% notation: #1

\hypertarget{prop-rolles-zero-derivative}{}
\begin{theorem}[Rolle's Theorem on the Zero of a Derivative]
\[
\TerForall \Function{f} \colon \ClosedInterval{0, 1} \TermTo \RealNumbers
\left( \Continuous{f} \And \TerForall a \in \OpenInterval{0, 1} \DerivativeAtPoint{f'}(a) = 0 \right)
\TermImplication f(0) = f(1)
\]
\end{theorem}

\begin{proof}[RU]
Пусть $\Function{g}(x) = \Function{f}(x) + cx$. Так как $\DerivativeAtPoint{g'}(x) = \DerivativeAtPoint{f'}(x) + c = c$, то из теоремы о среднем значении следует, что $\Function{f}(1) - \Function{f}(0) > -c$. Аналогично, рассматривая $\Function{h}(x) = cx - \Function{f}(x)$, получаем $\Function{f}(1) - \Function{f}(0) < c$. Поскольку это верно для любого $\RealNumbers{c} > 0$, заключаем, что $\Function{f}(0) = \Function{f}(1)$.
\end{proof}

\begin{proof}[EN]
Let $\Function{g}(x) = \Function{f}(x) + cx$. Since $\DerivativeAtPoint{g'}(x) = \DerivativeAtPoint{f'}(x) + c = c$, it follows from the Mean Value Theorem that $\Function{f}(1) - \Function{f}(0) > -c$. Similarly, by considering $\Function{h}(x) = cx - \Function{f}(x)$, we obtain $\Function{f}(1) - \Function{f}(0) < c$. Since this holds for any $\RealNumbers{c} > 0$, we conclude that $\Function{f}(0) = \Function{f}(1)$.
\end{proof}